
\documentclass[11pt,fleqn]{article} 
\usepackage[margin=0.8in, head=0.8in]{geometry} 
\usepackage{amsmath, amssymb, amsthm}
\usepackage{fancyhdr} 
\usepackage{palatino, url, multicol}
\usepackage{graphicx, pgfplots} 
\usepackage[all]{xy}
\usepackage{polynom} 
%\usepackage{pdfsync} %% I don't know why this messes up tabular column widths
\usepackage{enumerate}
\usepackage{framed}
\usepackage{setspace}
\usepackage{array,tikz}

\pgfplotsset{compat=1.6}

\pgfplotsset{soldot/.style={color=black,only marks,mark=*}} \pgfplotsset{holdot/.style={color=black,fill=white,only marks,mark=*}}


\pagestyle{fancy} 
\lfoot{}
\rfoot{\S 3.2 \& 3.3}

\begin{document}
\renewcommand{\headrulewidth}{0pt}
\newcommand{\blank}[1]{\rule{#1}{0.75pt}}
\newcommand{\bc}{\begin{center}}
\newcommand{\ec}{\end{center}}
\renewcommand{\d}{\displaystyle}

\vspace*{-0.7in}

%%%%%%%%%intro page
\begin{center}
  \large
  \sc{Last of 3.2 Trig Integrals and 3.3 Trig Substitution}\\
   
\end{center}

\begin{enumerate}
\item Trigonometric Identities (Fill in the set of Pythagorean Identities.)\\

\begin{tabular}{ccc}
Half-Angle/Double Angle && Sum and Difference\\
\hline
&&\\
$\sin^2(x)=\frac{1}{2} ( 1-\cos(2x))$&&$\sin(ax)\cos(bx)=\frac{1}{2} ( \sin((a-b)x)+\sin((a+b)x))$\\ &&\\

 $\cos^2(x)=\frac{1}{2} ( 1+\cos(2x))$&&$\sin(ax)\sin(bx)=\frac{1}{2} ( \cos((a-b)x)-\cos((a+b)x))$\\ &&\\
 
$\sin(2x)=2\sin(x)\cos(x)$ &&$\cos(ax)\cos(bx)=\frac{1}{2} ( \cos((a-b)x)+\cos((a+b)x))$\\ &&\\
Pythagorean \quad \hspace{.8 in} \quad&&\\
\hline
\end{tabular}
\vspace{1in}


\item Compare the following three integrals. Which integration strategy works on which one?\\

\begin{tabular}{ccc}
(a) $\d \int \cos(x) \: dx$ &\quad \hspace{.5in} (b) $\d \int \cos^2(x) \: dx$ &\quad \hspace{1in} (c) $\d \int \cos^3(x) \: dx$ \\
\end{tabular}
\newpage
\begin{center} Section 3.3 \end{center}
\item Compare the following three integrals. Which ones can a Calculus I student integrate?\\
	\begin{tabular}{ccc}	
	(a) $\displaystyle{\int \frac{1}{\sqrt{9-x^2}}\: dx}$& \quad \hspace{.5in} (b)$\displaystyle{\int \frac{x}{\sqrt{9-x^2}}\: dx}$ &\quad \hspace{1in} (c) $\displaystyle{\int  \frac{x^2}{\sqrt{9-x^2}}\: dx}$
	\end{tabular}
	\vfill
\item Summary: If {\Large{$\sqrt{a^2-x^2}$}} appears in an integrand (\textbf{and} other techniques do not work), then\\
\vspace{.2in}
\newpage
\item Compare the following integrals. Which one can a Calculus I student integrate?\\
\begin{tabular}{ccc}	
	(a) $\displaystyle{\int \frac{dx}{{9+x^2}}}$&
	\quad \hspace{.5in} (b) $\displaystyle{\int \frac{dx}{\sqrt{9+x^2}}}$&\quad \hspace{.5in} (c)
	$\displaystyle{\int  \frac{dx}{\sqrt{x^2-9}}\: dx}$
	\end{tabular}
	\vfill
\item Summary: 
	\begin{itemize}
	\item If {\Large{$\sqrt{a^2+x^2}$}} appears in an integrand (\textbf{and} other techniques do not work), then\\
	\vspace{.2in}
	\item If {\Large{$\sqrt{x^2-a^2}$}} appears in an integrand (\textbf{and} other techniques do not work), then\\
	\vspace{.2in}
	\end{itemize}
\newpage
\begin{center} EXTRA PRACTICE \end{center}
\item Evaluate 
	\begin{enumerate}
	\item $\displaystyle{\int \frac{dx}{(4+x^2)^2}  }$
	\vfill
	\item $\displaystyle{\int \frac{dx}{(x^2-9)^{3/2}}}$
	\vfill
	\end{enumerate}
	


\end{enumerate}
\end{document}

